\section{Conclusion}

In this work, we introduced JoyAI-Talker, a full-duplex speech dialogue system built upon the decoupled Duplex-Thinker-Talker architecture, which separates conversational state coordination, cognitive planning, and expressive speech generation into distinct modular boundaries. Overall, this work introduces several key technological advancements. First, JoyAI-Talker delivers robust foundation model capabilities that natively support advanced conversational voice agents, enabling complex multi-turn speech reasoning, tool-use, and interactive agentic behaviors directly from raw speech queries. Second, through a unified speech-text joint training pipeline spanning pretraining to preference optimization, the system significantly alleviates the common "cognitive degradation" bottleneck to a certain extent, largely preserving the model's fundamental reasoning, STEM, and logical capabilities in the textual domain while expanding into speech. Third, the Persona-Adaptive Empathetic Response (PAER) framework, structured as a hierarchical cognitive pipeline, integrates non-verbal user-attribute perception with Chain-of-Thought (CoT) reasoning to formulate context-adaptive, emotionally intelligent responses. Fourth, the Talker module utilizes a decoupled, instruction-controllable generation paradigm, translating the Thinker's semantic outputs and speaking instructions into speaker-consistent, expressive speech with precise, token-level controls over utterance-level expressions and localized paralinguistic events. Finally, the state-driven, plug-and-play Joy-Duplex framework  interleaves semantic and interaction state tokens directly into the streaming decoding path, acting as a non-intrusive gating engine for real-time turn control. Comprehensive evaluations demonstrate that JoyAI-Talker achieves highly competitive performance across a diverse range of foundational text-to-text (T2T) and speech-to-text (S2T) benchmarks, comparing favorably with state-of-the-art baselines. Notably, on the rigorous Full-Duplex-Bench v1.5, Joy-Duplex delivers a well-balanced interaction policy, achieving a high responsiveness of $0.88$ under user interruptions while keeping false triggers extremely low under background speech, representing a robust and stable paradigm for fluid, natural speech dialogue.
\section{Limitations}
While JoyAI-Talker achieves state-of-the-art performance in highly empathetic voice interaction and robust state-driven full-duplex control, several inherent limitations remain and motivate our ongoing research.

First, the model's perception is currently focused on speech signals and speaker attributes, lacking a dedicated capability for general audio event reasoning and environmental sound comprehension.  In real-world interaction scenarios, non-speech acoustic cues---such as background noise, sirens, alarms, and music---carry critical contextual information.  Future work will prioritize the development of a unified audio-speech encoder and aligned training corpora to enable holistic environmental audio understanding and robust scene-aware reasoning.

Second, the large-scale Reinforcement Learning (RL) training phase for JoyAI-Talker is still in progress.  Without a dedicated RL alignment stage, the model's performance on highly complex, multi-constraint instructions and voice-driven agent invocation can be relatively constrained.  Beyond conventional optimization of speech-text consistency, we plan to further explore the utility of RL.  Specifically, we aim to design a unified speech-text joint reinforcement learning strategy, investigating how RL can be effectively utilized to boost long-term speech reasoning stability, improve complex instruction-following, and refine the execution of voice-driven agentic workflows.

Finally, Joy-Duplex currently operates as a decoupled, state-driven modular framework.  While this design provides robust plug-and-play flexibility and low training overhead, a fully end-to-end dual-stream architecture can achieve favorable performance in minimizing turn-taking latency and realizing continuous, concurrent input-output synchronization.  In future work, we plan to further explore and compare these two paradigms to optimize the trade-offs between modular state-driven coordination and unified, time-aligned dual-stream sequence modeling along a shared temporal axis.

\section*{Authors (alphabetical order of family name)}



\begin{description}
    \item[Core Contributors:] Yinhao Bai, Jinming Chen, Yafeng Chen, Wei Deng, Boya Dong, Nan Duan, Yu Gu, Weisheng Han, Yankun Huang, Ming Ke, Hao Li, Jingdong Li, Xiangyu Liang, Ning Liu, Yuan Liu, Ji Miao, Jiaqi Wang, Qi Wang, Wenchao Wang, Yuxuan Wang, Zhenfang Wang, Zhangyu Xiao, Chao Xue, Hongfei Xue, Fan Yu, Tianyi Zhang, Yuan Zhang, Yuqi Zhang, Lin Zhu
    \item[Contributors:] Haoran Gu, Yiheng Jiang, Jun Wang, Ziteng Wang
\end{description}


